\chapter{变量与补充材料}

\section{变量编码表}

本文变量编码包括 AI 介入和视频内容两个部分。AI 介入形式编码为 AI 音频、AI 画面和混合；AI 介入程度编码为浅度、中度和深度；AI 技术质感编码为粗糙、普通和精良。视频主体编码为人、景色和物；视频主题编码为自然风光、历史人文、现代城市/生活、家国情怀和剧情/整活；内容导向编码为信息实用型和娱乐审美型；视频风格编码为纪实/写实、卡通/艺术化和抽象/超现实。连续变量包括视频时长和发布距采集日天数，建模时分别使用 $\log(\text{视频时长}+1)$ 和 $\log(\text{发布距采集日天数}+1)$。

完整编码数据和统计结果已随论文分析材料一并保存，主要包括原始编码表、描述统计结果和扩展实验结果。

\section{主要回归结果}

主模型以 $\log(\text{总互动量}+1)$ 为因变量，样本量为 722，自变量数为 18，$R^2=0.502$，调整 $R^2=0.490$。显著结果包括：AI 技术质感中的普通和精良类别均显著正向影响传播效果；视频风格中的抽象/超现实类别显著正向影响传播效果；log 视频时长显著负向影响传播效果；log 发布距采集日天数显著正向影响传播效果。

完整回归系数、标准误、置信区间和影响百分比见论文配套分析结果表中的“主回归与省份稳健”和“主回归具体影响”工作表。

\section{稳健性检验结果}

本文从替换因变量、调整时间口径、加入省份固定效应和处理极端值四个方向进行稳健性检验。替换因变量时，分别以 log 点赞量、log 评论量、log 收藏量和 log 转发量作为因变量重复估计模型；调整时间口径时，以 $\log(\text{日均互动量}+1)$ 作为辅助结果变量重复估计模型；加入省份固定效应用于控制地区资源、账号基础和运营条件等相对稳定差异；处理极端值时，采用剔除总互动量 Top 1\%、剔除总互动量 Top 5\% 和对 log 总互动量进行 1\%-99\% 缩尾处理。结果显示，AI 技术质感和抽象/超现实风格在多数检验中保持稳定。

完整结果见论文配套分析结果表中的“稳健性分因变量回归”“稳健性显著结果摘要”和“极端值稳健性核心变量”工作表。

\section{高互动识别与地区热力矩阵补充结果}

本文在高互动视频画像的基础上补充二分类 Logit 识别实验，并绘制地区互动表现与内容特征热力矩阵。高互动二分类 Logit 以是否进入总互动量前 25\% 作为因变量，主要用于观察内容变量对高互动组识别的补充价值。由于 AI 技术质感中的粗糙质感组未出现高互动样本，本文采用带 L2 惩罚项的 Logit 模型，以避免普通 Logit 中类别分离导致系数不稳定。地区热力矩阵则用于展示不同省份/地区在样本量、总互动量中位数、高互动占比、精良质感占比和抽象/超现实占比上的相对差异。

完整结果见论文配套分析结果表中的“高互动 Logit 核心变量”“Logit 模型指标”和“地区热力图数据”工作表。
